%    Copyright 2021 Silvan Zahno, HEI-Vs
%
% Licensed under the Apache License, Version 2.0 (the "License");
% you may not use this file except in compliance with the License.
% You may obtain a copy of the License at
%
% http://www.apache.org/licenses/LICENSE-2.0
%
% Unless required by applicable law or agreed to in writing, software
% distributed under the License is distributed on an "AS IS" BASIS,
% WITHOUT WARRANTIES OR CONDITIONS OF ANY KIND, either express or implied.
% See the License for the specific language governing permissions and
% limitations under the License.

% -----------------------------------------------------------------------------
% HEi-VS Bachelor - Guide to Thesis Document
%
% This documents showns you have to use the template
% -----------------------------------------------------------------------------

% Document settings
\documentclass[a4paper,11pt,fleqn]{book}
\usepackage[utf8]{inputenc}
\usepackage[T1]{fontenc}
\usepackage[german,french,english]{babel}

% -----------------------------------------------------------------------------
% Preamble
% -----------------------------------------------------------------------------
%% ------------------------------------------------------------------------------
% Thesis metadata
% ------------------------------------------------------------------------------

% Thesis info
\newcommand{\ThesisTitle}{[ThesisTitle]}
\newcommand{\ThesisSubject}{[ThesisSubject]}
\newcommand{\Orientation}{Systems Engineering}
\newcommand{\Major}{Infotronics}
\newcommand{\Year}{2021}
\newcommand{\Version}{v0.1}
\newcommand{\SubmissionDate}{DD.MM.YYYY}
\newcommand{\Keywords}{keyword1, keyword2, keyword3}

% Author
\newcommand{\AuthorFirstName}{[FirstName]}
\newcommand{\AuthorLastName}{[LastName]}
\newcommand{\AuthorEmail}{[user@domain.com]}
\newcommand{\AuthorDegree}{Bachelor}

\newcommand{\Author}{\AuthorFirstName\,\AuthorLastName}

% Advisor
\newcommand{\AdvisorFirstName}{Silvan}
\newcommand{\AdvisorLastName}{Zahno}
\newcommand{\AdvisorSchool}{HEI-Vs}
\newcommand{\AdvisorResearchUnit}{ISI}
\newcommand{\Advisor}{Prof. \AdvisorFirstName\,\AdvisorLastName}

% Main expert
\newcommand{\ExpertFirstName}{[FirstName]}
\newcommand{\ExpertLastName}{[LastName]}
\newcommand{\Expert}{\ExpertFirstName\,\ExpertLastName}
\newcommand{\ExpertLab}{[Lab/Company]}

% University
\newcommand{\University}{HES-SO//Valais Wallis}
\newcommand{\Dean}{Prof. Pierre Pompili}

% Location & Current Date
\newcommand{\Date}{\today}
\newcommand{\Place}{Sion}
\newcommand{\Country}{Switzerland}

% Logo
\providecommand{\logo}{}
\renewcommand{\logo}{latex-project.pdf}
\newcommand{\logosmall}{latex-project.pdf}

% ------------------------------------------------------------------------------
% Page Content
%
%   Titlepage        : tp          - creates a titlepage
%   Table of Content : toc         - creates a table of contents
%   Table of Figures : tof         - creates a table of figures
%   Table of Tables  : tot         - creates a table of tables
%   Table of Listings: tol         - creates a table of tables
%   Preface          : pre         - creates a frontmatter
%   Appendices       : apd         - creates a appendices
%   Glossary         : glo         - creates a glossary
%   Acronyms         : acr         - creates a acronym list
%   Bibliography     : bib         - creates a bilbiography list
%   Watermark        : w           - creates a watermark overall pages
%   Linenumbers      : ln          - adds linenumbering to the full document
%   Comment          : cmt         - Comments from the todo package \todo \missingfigure etc.
%\usepackage[tp, toc, tof, tot, tol, pre, apd, glo, acr, bib, w, ln, cmt]{optional} % all optionals active
\usepackage[tp, toc, tof, tot, tol, pre, apd, glo, acr, bib]{optional} % only optionals for final thesis active
%\usepackage[none]{optional} % no optionals active         % your project data
% Custom metadata
% ------------------------------------------------------------------------------

% Thesis info
\newcommand{\ThesisTitle}{Guide to Thesis}
\newcommand{\ThesisSubject}{How to use \latex for your thesis}
\newcommand{\Orientation}{Systems Engineering}
\newcommand{\Major}{Infotronics}
\newcommand{\Year}{2021}
\newcommand{\Version}{v1.0}
\newcommand{\SubmissionDate}{\today}

% Author
\newcommand{\AuthorFirstName}{Silvan}
\newcommand{\AuthorLastName}{Zahno}
\newcommand{\AuthorEmail}{silvan.zahno@hevs.ch}
\newcommand{\Author}{\AuthorFirstName\,\AuthorLastName}

% Advisor
\newcommand{\AdvisorFirstName}{Silvan}
\newcommand{\AdvisorLastName}{Zahno}
\newcommand{\Advisor}{\AdvisorFirstName\,\AdvisorLastName}

% Main expert
\newcommand{\ExpertFirstName}{Pierre}
\newcommand{\ExpertLastName}{Pompili}
\newcommand{\Expert}{\ExpertFirstName\,\ExpertLastName}

% University
\newcommand{\University}{HES-SO//Valais Wallis}

% Location & Current Date
\newcommand{\Date}{\today}
\newcommand{\Place}{Sion}
\newcommand{\Country}{Switzerland}

% ------------------------------------------------------------------------------
% Page Content
\usepackage[tp, toc, tof, apd, glo, acr, bib]{optional} % only optionals for final thesis active

% ------------------------------------------------------------------------------
% Template settings
% ------------------------------------------------------------------------------

% ------------------------------------------------------------------------------
% General tools
%
\usepackage{etoolbox}

% ------------------------------------------------------------------------------
% Page style
%
\usepackage[margin=3cm, left=3.5cm, right=3.5cm, twoside=true]{geometry}
\usepackage{fancyhdr}
\setlength{\headheight}{14pt}
\renewcommand{\sectionmark}[1]{\markright{\thesection\ #1}}
\pagestyle{fancy}

% Standard pages (inside chapters)
\fancyhf{}
\usepackage{lastpage}
\renewcommand{\headrulewidth}{0.4pt}
\renewcommand{\footrulewidth}{0pt}
\fancyhead[OR]{\bfseries \nouppercase{\rightmark}}
\fancyhead[EL]{\bfseries \nouppercase{\leftmark}}
\fancyfoot[EL,OR]{\thepage/\pageref{LastPage}}

% First page of chapters
\fancypagestyle{plain}{
  \fancyhf{}
  \renewcommand{\headrulewidth}{0pt}
  \renewcommand{\footrulewidth}{0pt}
  \fancyfoot[EL,OR]{\thepage/\pageref{LastPage}}
}

% Imports for external PDFs
\fancypagestyle{addpagenumbersforpdfimports}{
  \fancyhead{}
  \renewcommand{\headrulewidth}{0pt}
  \fancyfoot{}
  \fancyfoot[RO,LE]{\thepage/\pageref{LastPage}}
}

% Use empty style for page when clearing double pages
\def\cleartoodd{%
  \clearpage%
  \ifodd\value{page}\else\mbox{}\thispagestyle{empty}\newpage\fi%
}

\def\clearchap{%
  \ifodd\value{page}\else\mbox{}\thispagestyle{empty}\fi%
}

% \cleardoublepage replaced by \cleartoodd
\let\origdoublepage\cleardoublepage
\renewcommand{\cleardoublepage}{%
  \cleartoodd%
}

% ------------------------------------------------------------------------------
% Colors
%
\usepackage{xcolor}    % colors
% HESSO VS Colors
\definecolor{HEICyan}{HTML}{2699D4}
\definecolor{HEIMagenta}{HTML}{C8306F}
\definecolor{HEIYellow}{HTML}{E19A27}
\definecolor{HEIGreen}{HTML}{238B56}

% Smart Process Lab Colors
\definecolor{SPLGreen}{RGB}{190, 214, 0}
\definecolor{SPLBlue}{RGB}{0, 169, 224}
\definecolor{SPLPurple}{HTML}{C8306F}

% Standard Colors
\definecolor{mGray80}{gray}{0.8}
\definecolor{mGray60}{gray}{0.6}
\definecolor{mGray40}{gray}{0.4}
\definecolor{mGray20}{gray}{0.2}
\definecolor{mWhite}{RGB}{255, 255, 255}
\definecolor{mBlack}{RGB}{0, 0, 0}
\definecolor{mPink}{cmyk}{0, 0.7808, 0.4429, 0.1412}

% Document color accents
\definecolor{coloraccent1}{HTML}{C8306F}
\definecolor{coloraccent2}{HTML}{2699D4}
\definecolor{coloraccent3}{HTML}{238B56}
\definecolor{coloraccent4}{HTML}{E19A27}

% Hyperlink colors
\definecolor{linkColor}{HTML}{2C89A0}
\definecolor{citeColor}{HTML}{2C89A0}
\definecolor{fileColor}{HTML}{2C89A0}
\definecolor{urlColor}{HTML}{2C89A0}
\definecolor{linkColor}{HTML}{2C89A0}

% Todo Colors
\definecolor{todoColor}{RGB}{3, 169, 244}
\definecolor{unsureColor}{RGB}{255, 165, 0}
\definecolor{infoColor}{RGB}{139, 195, 74}
\definecolor{criticalColor}{RGB}{211, 47, 47}
\definecolor{changeColor}{RGB}{142, 68, 173}
\definecolor{questionColor}{RGB}{156, 39, 176}
\definecolor{improvementColor}{RGB}{230, 126, 34}

% ------------------------------------------------------------------------------
% Fonts
%
% Helvetica
\usepackage{helvet} % helvetica
\usepackage{pifont} % icon package <http://willbenton.com/wb-images/pifont.pdf>

% ------------------------------------------------------------------------------
% Special characters
%
\newcommand{\q}{\textnormal{\textquotesingle}}
\newcommand{\ql}{\textnormal{\textquoteleft}}
\newcommand{\qr}{\textnormal{\textquoteright}}
\newcommand{\qq}{\textnormal{\textquotedblright}}
\newcommand{\qql}{\textnormal{\textquotedblleft}}
\newcommand{\qqr}{\textnormal{\textquotedblright}}

% ------------------------------------------------------------------------------
% Todo List
% <https://tex.stackexchange.com/questions/247681/how-to-create-checkbox-todo-list>
\usepackage{enumitem}
\newlist{todolist}{itemize}{2}
\setlist[todolist]{label=$\square$}
\setlist[itemize]{noitemsep, nolistsep}
\setlist[enumerate]{noitemsep, nolistsep}
\newcommand{\cmark}{\ding{51}}
\newcommand{\xmark}{\ding{55}}
\newcommand{\done}{\rlap{$\square$}{\raisebox{2pt}{\large\hspace{1pt}\cmark}} \hspace{-2.5pt}}
\newcommand{\wontfix}{\rlap{$\square$}{\large\hspace{1pt}\xmark}}

% Custom todo notes
\usepackage{todonotes}
\opt{cmt}{
  \setuptodonotes{backgroundcolor=todoColor!50,bordercolor=todoColor!75,linecolor=todoColor!75,size=\scriptsize}
}
\setuptodonotes{backgroundcolor=todoColor!50,bordercolor=todoColor!75,linecolor=todoColor!75,size=\scriptsize,disable}
%\usepackage[backgroundcolor=todoColor!50,bordercolor=todoColor!75,linecolor=todoColor!75,size=\scriptsize,disable]{todonotes}
\newcommand{\unsure}[1]{\todo[backgroundcolor=unsureColor!50,bordercolor=unsureColor!75,linecolor=unsureColor!75,size=\scriptsize]{#1}}
\newcommand{\info}[1]{\todo[backgroundcolor=infoColor!50,bordercolor=infoColor!75,linecolor=infoColor!75,size=\scriptsize]{#1}}
\newcommand{\critical}[1]{\todo[color=criticalColor!50,bordercolor=criticalColor!75,linecolor=criticalColor!75,size=\scriptsize]{#1}}
\newcommand{\change}[1]{\todo[color=changeColor!50,bordercolor=changeColor!75,linecolor=changeColor!75,size=\scriptsize]{#1}}
\newcommand{\question}[1]{\todo[color=questionColor!50,bordercolor=questionColor!75,linecolor=questionColor!75,size=\scriptsize]{#1}}
\newcommand{\improvement}[1]{\todo[backgroundcolor=improvementColor!50,bordercolor=improvementColor!75,linecolor=improvementColor!75,size=\scriptsize]{#1}}

% ------------------------------------------------------------------------------
% Math
%
\usepackage{amsmath}  % better math
\usepackage{amssymb}

% ------------------------------------------------------------------------------
% Floats and figures
%
\usepackage{newfloat}          % floats
\usepackage[twoside]{caption}  % captions
\usepackage{subcaption}        % subcaptions
\usepackage[section]{placeins} % allows to put float barriers

% Float captions in italics, with label in margin
\DeclareCaptionLabelFormat{title}{#1 #2}
\DeclareCaptionLabelFormat{hangout}{\llap{#1 #2\hspace{5mm}}}
\captionsetup{
  format=hang,
  labelformat=hangout,
  singlelinecheck=false,
  font={it}
}

% ------------------------------------------------------------------------------
% Tables
%
\usepackage{arydshln}  % dashed vertical and horizontal lines
\usepackage{booktabs}  % Thick horizontal lines
\usepackage{hhline}    % custom horizontal lines
\usepackage{boldline}  % Bold tables lines
\usepackage{multirow}  % vertically merged rows
\usepackage{slashbox}  % topleft cell can be split
\usepackage{longtable} % Table span across multiple pages
\usepackage{array}     % manipulate array
%\usepackage{tabularx}  % better tables

% ------------------------------------------------------------------------------
% Page Overlays
%
\usepackage{lineno}                      % Line numbering
\opt{ln}{
  \linenumbers
}

\opt{w}{
  \usepackage{tikz}
  \usepackage[printwatermark]{xwatermark}  % Watermark
  \newcommand{\watermarkText}{Confidential \& Draft \& \Version}
  \newcommand{\watermarkOpacity}{0.3}
  \newcommand{\watermarkColor}{red}
  \newsavebox\mybox
  \savebox\mybox{\tikz[color=\watermarkColor,opacity=\watermarkOpacity]\node{\watermarkText};}
  \newwatermark*[
    allpages,
    angle=45,
    scale=4,
    xpos=-20,
    ypos=0
  ]{\usebox\mybox}
}

% ------------------------------------------------------------------------------
% Bibliography
%
\opt{bib}{
  % Use biber, with numeric style and no sorting (citation order)
  \usepackage[backend=biber,style=numeric,sorting=none,bibencoding=auto]{biblatex}
  \addbibresource{03-tail/bibliography.bib}
}

% ------------------------------------------------------------------------------
% Tables of contents, figures, tables and listings
%
\usepackage{tocloft}
\newlistof{listing}{lol}{List of Listings}
\setcounter{tocdepth}{1} % Depth to 'section'
\setlength{\cftfigindent}{0pt}  % remove indentation from figures in lof
\setlength{\cftfignumwidth}{1cm}
\setlength{\cfttabindent}{0pt}  % remove indentation from tables in lot
\setlength{\cfttabnumwidth}{1cm}
\setlength{\cftlistingindent}{0pt}
\setlength{\cftlistingnumwidth}{1cm}

% ------------------------------------------------------------------------------
% Mini tables of contents
%
\usepackage{minitoc}

% no "Contents" title
\mtcsettitle{minitoc}{Contents}

% Layout
\setlength{\mtcindent}{-0.5em}
\mtcsetoffset{minitoc}{-1em}

% Spacing above and below table
\mtcsetfeature{minitoc}{before}{\vspace{0.5cm}}
\mtcsetfeature{minitoc}{after}{\vspace{-0.25cm}}
\renewcommand{\mtifont}{\sffamily\bfseries\large}

% ------------------------------------------------------------------------------
% Graphics
%
\usepackage[xetex]{graphicx} % graphics importing
\graphicspath{
  {02-main/figures/}
  {04-resources/}
}

% ------------------------------------------------------------------------------
% Code and syntax highlighting
%
\usepackage[newfloat]{minted}   % code highlighting
\setminted{frame=single, framesep=6pt, breaklines=true, fontsize=\scriptsize}
% all - manni, fruity, rrt, autumn, perldoc, bw, borland, emacs, colorful,
% vim murphy, pastie, vs, friendly, trac, native, tango, monokai
\newcommand{\codecolorscheme}{pastie}       % bright good for printing
% \newcommand{\codecolorscheme}{monokai}    % dark good for viewing needs dark background
% \newcommand{\codecolorscheme}{bw}         % black and white good with light gray background
\newcommand{\codecolorschemebg}{gray!20}    % bright background good for printing
% \newcommand{\codecolorschemebg}{gray!80}  % dark background bright good for viewing

% ------------------------------------------------------------------------------
% Typography
%
\usepackage{csquotes}                    % paragraph indentation and spacing
\usepackage[defaultlines=3,all]{nowidow} % avoid widows and orphans
%\usepackage{microtype}                   % typographic improvements
\usepackage{parskip}                     % No indent and auto-space between paragraphs
%\usepackage{fontspec}                    % To change fonts in XeTex
\usepackage[super]{nth}

\usepackage{paralist}
\setlist{after=\vspace{\baselineskip}}

% ------------------------------------------------------------------------------
% Section and chapters headings
%
\usepackage[explicit]{titlesec} % titles formatting
%\usepackage{titletoc} % titles formatting in ToC etc
%\usepackage{sectsty}  % sectioning commands

% -- Chapters --
% Remove "Chapter N" and use a sans-serif font

% Set layout lengths
\setlength{\headheight}{8mm}
\setlength{\footskip}{1.5cm}
\addtolength{\textheight}{-.5cm}

\titlespacing{\chapter}{-5mm}{-10mm}{3mm}
\titlespacing{\section}{-5mm}{3mm}{2mm}
\titlespacing{\subsection}{-5mm}{2mm}{2mm}
\titlespacing{\subsubsection}{-5mm}{2mm}{2mm}

%\titleformat{\chapter}[block]
%{\Huge}
%{\thechapter\hspace{12pt}\textcolor{gray80}{|}\hspace{12pt}}
%{0pt}
%{\Huge\bfseries}

\titleformat{\chapter}{\Huge\bfseries}{\llap{\thechapter\hspace{12pt}\textcolor{mGray80}{|}}}{0mm}{%
  \hfill\begin{minipage}[t]{\dimexpr\textwidth}\raggedright#1\end{minipage}%
}
\titleformat{\section}{\Large\bfseries}{\llap{\thesection}}{0mm}{%
  \hfill\begin{minipage}[t]{\dimexpr\textwidth}\raggedright#1\end{minipage}%
}
\titleformat{\subsection}{\large \bfseries}{\llap{\thesubsection}}{0mm}{%
  \hfill\begin{minipage}[t]{\dimexpr\textwidth}\raggedright#1\end{minipage}%
}
\titleformat{\subsubsection}{\bfseries}{\llap{\thesubsubsection}}{0mm}{%
  \hfill\begin{minipage}[t]{\dimexpr\textwidth}\raggedright#1\end{minipage}%
}

% ------------------------------------------------------------------------------
% Misc
%
\usepackage{lipsum}     % filler text
\usepackage{blindtext}  % random text
\usepackage{lscape}     % easy landscape pages
\usepackage{pdflscape}  % landscape pages for PDFs
\usepackage[final]{pdfpages} %
\usepackage{nameref}    % name referencing
\usepackage{awesomebox} % info, warning, tip Boxes
\usepackage{multicol}   % multicols environment

% ------------------------------------------------------------------------------
% References
%
\usepackage{url}

% pdf metadata
\usepackage[
  pdfauthor={\Author},
  pdftitle={\ThesisTitle},
  pdfsubject={\ThesisSubject},
  pdfkeywords={\Keywords}
  pdfduplex=DuplexFlipLongEdge]{hyperref}

% Hyperlinks
\hypersetup{
  colorlinks=true,
  linkcolor=linkColor,
  citecolor=citeColor,
  filecolor=fileColor,
  urlcolor=urlColor,
  pdfborder= 0.1 0.1 1
}
\opt{nlb}{
    \hypersetup{pdfborder = 0 0 0} % remove colored border for links in pdf
}

% ------------------------------------------------------------------------------
% Glossary
%
\opt{acr,glo}{
  \usepackage[acronym,toc,section=chapter]{glossaries}
  %\usepackage[xindy,toc]{glossaries}
  % ------------------------------------------------------------------------------
% Terms
% -----
% format:  \newglossaryentry{<label>}{<settings>}
% example: \newglossaryentry{computer}
%{
%	name=computer,
%	description={is a programmable machine that receives input, stores and manipulates data, and provides output in a useful format}
%}
% Development
% ------------------------------------------------------------------------------

\newglossaryentry{scrum}{
    name=Scrum,
    description={Scrum is an agile process framework for managing complex knowledge work, with an initial emphasis on software development, although it has been used in other fields and is slowly starting to be explored for other complex work, research and advanced technologies.}
}

% ------------------------------------------------------------------------------
% Acronyms
% --------
% format:  \newacronym{<label>}{<abbrv>}{<full>}
% example: \newacronym{lvm}{LVM}{Logical Volume Manager}
% plural:  \newacronym[longplural={Frames per Second}]{fpsLabel}{FPS}{Frame per Second}
% ------------------------------------------------------------------------------

% Machine Learning
\newacronym{ar}{AR}{Augmented Reality}

% Else
\newacronym{uv}{UV}{Ultra Violet}


  \makeglossaries
}
\usepackage{amsmath}  % better math
\usepackage{amssymb}    % template settings
% ------------------------------------------------------------------------------
% Custom Template settings can be overwritten in this document
% ------------------------------------------------------------------------------

\usepackage{cmbright}       % Sans serif
\usepackage{sourcecodepro}  % Monospace
\renewcommand{\familydefault}{\sfdefault}  % your custom packages etc

\begin{document}
% -----------------------------------------------------------------------------
% Front matter
% -----------------------------------------------------------------------------
\frontmatter

\dominitoc
% ------------------------------------------------------------------------------
% HES-SO Valais Wallis thesis title page
% ------------------------------------------------------------------------------

\begin{titlepage}
\newgeometry{margin=2.5cm}
 \scalefont{2}

{%\fontfamily{phv}\fontseries{mc}\selectfont
  \begin{center}
    \begin{minipage}{0.5\textwidth}
      \begin{flushleft}
        \includegraphics[width=0.75\columnwidth]{logos/hei-vs-light.pdf}
      \end{flushleft}
    \end{minipage}%
    \begin{minipage}{0.5\textwidth}
      \begin{flushright}
        \includegraphics[width=0.6\columnwidth]{logos/hesso-logo.pdf}
      \end{flushright}
    \end{minipage}

    \vspace{1.0cm}
    {
    \huge
    Degree Programme\\[0.5cm]
    \Orientation\\[0.5cm]
    }

    {
    \Large
    Major \Major \\[1.0cm]
    }

    {
    \Huge
    \textbf{
      BACHELOR'S THESIS\\[1.0cm]
      DIPLOMA \Year\\[1.5cm]
    }
    }

    {
    \huge
    \Author\\[1.7cm]
    }

    % Title
    {
      \Huge
      \textit{\ThesisTitle}
    }
  \end{center}
  \begin{picture}(0, 0)
    \linethickness{1 pt}
    \put(0, 70){\line(1, 0){460}}
    \put(0, 8){\line(0, 1){62}}
    \linethickness{2 pt}
    \put(0, 8){\line(1, 0){460}}
    \put(460, 8){\line(0, 1){62}}
  \end{picture}
   \vspace{1.0cm}

  \begin{flushleft}
    {
    \normalsize
    \textit{Professor} \\
    \Advisor \\[0.2cm]
    \textit{Expert} \\
    \Expert \\[0.2cm]
    \textit{Submission date of the report} \\
    \SubmissionDate \\[1.0cm]
    }
  \end{flushleft}

  \vfill
  \begin{minipage}{0.5\textwidth}
    \begin{flushleft}
      \includegraphics[width=0.5\columnwidth]{logos/hevs-pictogram-only.pdf}
    \end{flushleft}
  \end{minipage}%
  \begin{minipage}{0.5\textwidth}
    \begin{flushright}
      \includegraphics[width=0.2\columnwidth]{logos/valais-excellence-logo.pdf}
    \end{flushright}
  \end{minipage}

\scalefont{0.5}
\vspace{5pt}

}
\restoregeometry
\end{titlepage}





% Table of contents
\opt{toc}{
  \newpage
  \phantomsection
  \addcontentsline{toc}{chapter}{Contents}
  \tableofcontents
}

% Restore paragraphs
\setlength{\parskip}{1em}

% -----------------------------------------------------------------------------
% Main matter
% -----------------------------------------------------------------------------
\mainmatter

\setcounter{mtc}{1} % Help minitoc skip the front matter chapters
% -----------------------------------------------------------------------------
\chapter{Installation}
\minitoc\newpage

\begin{enumerate}
  \item Install \href{MiKTeX}{https://miktex.org/download}
  \item MiKTeX Packages
  \begin{itemize}
    \item minted
    \begin{itemize}
      \item pip install pygments
      \item add Python Scripts to \texttt{PATH} Environment Variable
    \end{itemize}
  \end{itemize}
  \item Install your preferred IDE
  \begin{itemize}
    \item We use \href{https://code.visualstudio.com/}{Visual Studio Code}
    \item With the following extensions
    \begin{itemize}
      \item \href{https://marketplace.visualstudio.com/items?itemName=torn4dom4n.latex-support}{LaTeX language support}
      \item \href{https://marketplace.visualstudio.com/items?itemName=James-Yu.latex-workshop}{LaTeX Workshop}
    \end{itemize}
  \end{itemize}
\end{itemize}

% -----------------------------------------------------------------------------
\chapter{Page Formatting}
\minitoc\newpage

\section{General Formatting}
\subsection{Page Formatting}
\begin{minted}{latex}
\newpage        % jump to new page

\par            % new paragraph
\\              % new paragraph
\\* or \newline % new line

%               % start line comment


\end{minted}

\subsection{Font sizes}
\begin{multicols}{2}
  {\Huge Huge}\\{\huge huge}\\{\LARGE LARGE}\\{\Large Large}\\{\large large}\\{\normalsize normal (default)}\\{\small small}\\{\footnotesize footnotesize}\\{\scriptsize scriptsize}\\{\tiny tiny}

  \vfill\null\columnbreak

  \begin{minted}{latex}
{\Huge Huge}
{\huge huge}
{\LARGE LARGE}
{\Large Large}
{\large large}
{\normalsize normal (default)}
{\small small}
{\footnotesize footnotesize}
{\scriptsize scriptsize}
{\tiny tiny}
  \end{minted}
\end{multicols}

\subsection{Decorations}
\begin{multicols}{2}
  \textit{Italic}\\\texttt{Typewriter}\\\textbf{Bold}\\\texttt{Text}\\\textsf{Serif Font}\\\textrm{Serif (Roman)}\\\underline{Underline}\\\emph{Emphasis}\\

  \vfill\null\columnbreak

  \begin{minted}{latex}
\textit{Italic}
\texttt{Typewriter}
\textbf{Bold}
\texttt{Text}
\textsf{Serif}
\textrm{Sans Serif (Roman)}
\underline{Underline}
\emph{Emphasis}
  \end{minted}
\end{multicols}

\subsection{Text Alignment}
\begin{multicols}{2}
  \raggedright
  left aligned text \\
  \raggedleft
  right aligned text \\
  \centering
  centered text \\
  %\justify
  %justified text \\

  \vfill\null\columnbreak

  \begin{minted}{latex}
\raggedright
left aligned text \\
\raggedleft
right aligned text \\
\centering
centered text \\
%\justify
%justified text \\
  \end{minted}
\end{multicols}

\subsection{Links}
\begin{multicols}{2}
  \url{http://hevs.ch} \\
  \href{mailto:\AuthorEmail}{\AuthorEmail} \\
  \href{http://hevs.ch}{Hevs Website}

  \vfill\null\columnbreak

  \begin{minted}{latex}
\url{http://hevs.ch}
\href{mailto:\AuthorEmail}{\AuthorEmail}
\href{http://hevs.ch}{Hevs Website}
  \end{minted}
\end{multicols}

\subsection{Cross referencing} \label{sec:crossref}
Set labels and reference them afterwards. Labels can be set anywhere hereafter examples for sections, equations and images.

\ref{sec:crossref} \nameref{sec:crossref}

\begin{equation} \label{eq:1}
  \sum_{i=0}^{\infty} a_i x^i
\end{equation}
The equations \ref{eq:1} is a power series.

\begin{center}
  \begingroup
  \includegraphics[width=0.5\columnwidth]{placeholder.pdf}
  \captionof{figure}{Figure with reference label}
  \label{fig:figurelabel}
  \endgroup
\end{center}
Figure \ref{fig:figurelabel} shows the image.

\begin{minted}{latex}
\section{Cross referencing} \label{crossref}
\ref{crossref} \nameref{crossref}

\begin{equation} \label{eq:1}
  \sum_{i=0}^{\infty} a_i x^i
\end{equation}
The equations \ref{eq:1} is a power series.

\begin{center}
  \begingroup
  \includegraphics[width=0.5\columnwidth]{placeholder.pdf}
  \captionof{figure}{Figure with reference label}
  \label{fig:figurelabel}
  \endgroup
\end{center}
figure \ref{fig:figurelabel} shows the logo.
\end{minted}

% -----------------------------------------------------------------------------
\chapter{General Elements}
\minitoc\newpage

\section{Sections}
For creating Chapters, Sections and Subsections there are multiple levels available.
Good practice is to limit the number of level to \mintinline{latex}{\subsubsection}
\begin{minted}{latex}
 0 = \chapter{Chapter}
 1 = \section{First Section}
 2 = \subsection{Second Section}
 3 = \subsubsection{Third Section}
 4 = \paragraph{Paragraph}
 5 = \subparagraph{Subpararaph}
\end{minted}

\section{Lists}
\begin{multicols}{2}
  \begin{itemize}
    \item One
    \item Two
  \end{itemize}

  \vfill\null\columnbreak

  \begin{minted}{latex}
\begin{itemize}
  \item One
  \item Two
\end{itemize}
   \end{minted}
\end{multicols}

\begin{multicols}{2}
  \begin{itemize}
    \item One
    \begin{itemize}
      \item Two
      \item Three
    \end{itemize}
    \item Four
  \end{itemize}

  \vfill\null\columnbreak

  \begin{minted}{latex}
\begin{itemize}
  \item One
  \begin{itemize}
    \item Two
    \item Three
  \end{itemize}
  \item Four
\end{itemize}
  \end{minted}
\end{multicols}

\begin{multicols}{2}
  \begin{enumerate}
    \item One
    \item Two
  \end{enumerate}

  \vfill\null\columnbreak

  \begin{minted}{latex}
\begin{enumerate}
  \item One
  \item Two
\end{enumerate}
  \end{minted}
\end{multicols}

\begin{multicols}{2}
  \begin{enumerate}
    \item One
    \begin{enumerate}
      \item Two
      \item Three
    \end{enumerate}
    \item Four
  \end{enumerate}

  \vfill\null\columnbreak

  \begin{minted}{latex}
\begin{enumerate}
  \item One
  \begin{enumerate}
    \item Two
    \item Three
  \end{enumerate}
  \item Four
\end{enumerate}
  \end{minted}
\end{multicols}

\begin{multicols}{2}
  \begin{todolist}[noitemsep]
    \item Normal item
    \item[\cmark] cmark item
    \item[\xmark] xmark item
    \item[\done] done item
    \item[\wontfix] wontfix item
  \end{todolist}

  \vfill\null\columnbreak

  \begin{minted}{latex}
\begin{todolist}
  \item Normal item
  \item[\cmark] cmark item
  \item[\xmark] xmark item
  \item[\done] done item
  \item[\wontfix] wontfix item
\end{todolist}
  \end{minted}
\end{multicols}

\section{Multicolumns}
\begin{multicols}{2}
  \begin{multicols}{2}
    Column 1
    \vfill\null\columnbreak
    Column 2
  \end{multicols}

  \vfill\null\columnbreak

  \begin{minted}{latex}
\begin{multicols}{2}
  Column 1
  \vfill\null\columnbreak
  Column 2
\end{multicols}
  \end{minted}
\end{multicols}

% -----------------------------------------------------------------------------
\chapter{Images}
\minitoc\newpage

\begin{multicols}{2}
  \begin{center}
    \begingroup
    \includegraphics[width=0.5\columnwidth]{placeholder.pdf}
    \captionof{figure}{Defined figure placement}
    \label{fig:figureplacement}
    \endgroup
  \end{center}

  \vfill\null\columnbreak

  \begin{minted}{latex}
\begin{center}
  \begingroup
  \includegraphics[width=0.5\columnwidth]placeholder.pdf}
  \captionof{figure}{Defined figure placement}
  \label{fig:figureplacement}
  \endgroup
\end{center}
  \end{minted}
\end{multicols}

\begin{multicols}{2}
  \begin{figure}
    \includegraphicswidth=0.5\columnwidth]{placeholder.pdf}
    \captionof{figure}{Automatic figure placement}
    \label{fig:automaticfigureplacement}
  \end{figure}

  \vfill\null\columnbreak
  \begin{minted}{latex}
\begin{figure}
  \includegraphicswidth=0.5\columnwidth]{placeholder.pdf}
  \captionof{figure}{Automatic figure placement}
  \label{fig:automaticfigureplacement}
\end{figure}
  \end{minted}
\end{multicols}

\begin{multicols){2}

  This section contains examples of figures: \ref{fig:subimages}, \ref{fig:topleft}, \ref{fig:topright}, \ref{fig:bottomleft}, \ref{fig:bottomright}

  \begin{figure}
    \centering
    \begin{subfigure}[b]{0.45\textwidth}
      \includegraphics[width=\textwidth]{placeholder.pdf}
      \caption{topleft}
      \label{fig:topleft}
    \end{subfigure}
    ~ %add desired spacing between images, e. g. ~, \quad, \qquad, \hfill etc.
    %(or a blank line to force the subfigure onto a new line)
    \begin{subfigure}[b]{0.45\textwidth}
      \includegraphics[width=\textwidth]{placeholder.pdf}
      \caption{topright}
      \label{fig:topright}
    \end{subfigure}
    ~ %add desired spacing between images, e. g. ~, \quad, \qquad, \hfill etc.
    %(or a blank line to force the subfigure onto a new line)
    \begin{subfigure}{0.45\textwidth}
      \includegraphics[width=\textwidth]{placeholder.pdf}
      \caption{bottomleft}
      \label{fig:bottomleft}
    \end{subfigure}
    ~ %add desired spacing between images, e. g. ~, \quad, \qquad, \hfill etc.
    %(or a blank line to force the subfigure onto a new line)
    \begin{subfigure}{0.45\textwidth}
      \includegraphics[width=\textwidth]{placeholder.pdf}
      \caption{bottomright}
      \label{fig:bottomright}
    \end{subfigure}
    \caption{Example subcaption}\label{fig:subimages}
  \end{figure}

  \vfill\null\columnbreak
  \begin{minted}{latex}
This section contains examples of figures: \ref{fig:floatingimage}, \ref{fig:topleft}, \ref{fig:topright}, \ref{fig:bottomleft}, \ref{fig:bottomright}

\begin{figure}
  \centering
  \begin{subfigure}[b]{0.45\textwidth}
    \includegraphics[width=\textwidth]{placeholder.pdf}
    \caption{topleft}
    \label{fig:topleft}
  \end{subfigure}
  ~ %add desired spacing between images, e. g. ~, \quad, \qquad, \hfill etc.
  %(or a blank line to force the subfigure onto a new line)
  \begin{subfigure}[b]{0.45\textwidth}
    \includegraphics[width=\textwidth]{placeholder.pdf}
    \caption{topright}
    \label{fig:topright}
  \end{subfigure}
  ~ %add desired spacing between images, e. g. ~, \quad, \qquad, \hfill etc.
  %(or a blank line to force the subfigure onto a new line)
  \begin{subfigure}{0.45\textwidth}
    \includegraphics[width=\textwidth]{placeholder.pdf}
    \caption{bottomleft}
    \label{fig:bottomleft}
  \end{subfigure}
  ~ %add desired spacing between images, e. g. ~, \quad, \qquad, \hfill etc.
  %(or a blank line to force the subfigure onto a new line)
  \begin{subfigure}{0.45\textwidth}
    \includegraphics[width=\textwidth]{placeholder.pdf}
    \caption{bottomright}
    \label{fig:bottomright}
  \end{subfigure}
  \caption{Example subcaption}\label{fig:subimages}
\end{figure}
  \end{minted}
\end{multicols}

% -----------------------------------------------------------------------------
\chapter{Tables}
\minitoc\newpage

\begin{multicols}{2}
  \centering
  \begingroup
  \begin{tabular}{l : c | r || r V{2.7} r \vline}
    Left & Center & Right & Right & Right \\
    1.1 & 1.2 & 1.3 & 1.4 & 1.5 \\ \hdashline
    2.1 & 2.2 & 2.3 & 2.4 & 2.5 \\ \hline
    3.1 & 3.2 & 3.3 & 3.4 & 3.5 \\ \hline \hline
    4.1 & 4.2 & 4.3 & 4.4 & 4.5 \\ \hlineB{2.7}
    5.1 & 5.2 & 5.3 & 5.4 & 5.5 \\ \toprule
    6.1 & 6.2 & 6.3 & 6.4 & 6.5 \\ \midrule
    7.1 & 7.2 & 7.3 & 7.4 & 7.5 \\ \bottomrule
  \end{tabular}
  \captionof{table}{Table Example 1}
  \label{tab:ex1}
  \endgroup

  \vfill\null\columnbreak
  \centering
  \begin{minted}{latex}
\begingroup
\begin{tabular}{l : c | r || r V{2.7} r \vline}
  Left & Center & Right & Right & Right \\
  1.1 & 1.2 & 1.3 & 1.4 & 1.5 \\ \hdashline
  2.1 & 2.2 & 2.3 & 2.4 & 2.5 \\ \hline
  3.1 & 3.2 & 3.3 & 3.4 & 3.5 \\ \hline \hline
  4.1 & 4.2 & 4.3 & 4.4 & 4.5 \\ \hlineB{2.7}
  5.1 & 5.2 & 5.3 & 5.4 & 5.5 \\ \toprule
  6.1 & 6.2 & 6.3 & 6.4 & 6.5 \\ \midrule
  7.1 & 7.2 & 7.3 & 7.4 & 7.5 \\ \bottomrule
\end{tabular}
\captionof{table}{Table Example 1}
\label{tab:ex1}
\endgroup

   \end{minted}
\end{multicols}

\begin{multicols}{2}
  \begingroup
  \begin{tabular}{c V{2.7} c  c V{2.7} c}
    X & $Q_1$ & $Q_0$ & Y \\ \hline \hline
    0 & 0 & 0 &  0 \\ \hline
    0 & 0 & 1 &  1 \\ \hline
    0 & 1 & 0 &  1 \\ \hline
    0 & 1 & 1 &  0 \\ \hlineB{2.7}
    1 & 0 & 0 &  1 \\ \hline
    1 & 0 & 1 &  1 \\ \hline
    1 & 1 & 0 &  1 \\ \hline
    1 & 1 & 1 &  0 \\ \hline
  \end{tabular}
  \captionof{table}{Table Example 2}
  \label{tab:ex2}
  \endgroup

  \vfill\null\columnbreak

  \begin{minted}{latex}
\begingroup
\begin{tabular}{c V{2.7} c  c V{2.7} c}
  X & $Q_1$ & $Q_0$ & Y \\ \hline \hline
  0 & 0 & 0 &  0 \\ \hline
  0 & 0 & 1 &  1 \\ \hline
  0 & 1 & 0 &  1 \\ \hline
  0 & 1 & 1 &  0 \\ \hlineB{2.7}
  1 & 0 & 0 &  1 \\ \hline
  1 & 0 & 1 &  1 \\ \hline
  1 & 1 & 0 &  1 \\ \hline
  1 & 1 & 1 &  0 \\ \hline
\end{tabular}
\captionof{table}{Table Example 2}
\label{tab:ex2}
\endgroup

  \end{minted}
\end{multicols}

\begin{multicols}{2}
  \begingroup
  \begin{tabular}{ | c | c | }
    \hline
    \multicolumn{1}{|c|}{\emph{ Operator }} & \multicolumn{1}{|c|}{\emph{ Beschreibung}} \\ \hline
    +  & \multicolumn{1}{|l|}{Addition} \\ \hline
    -  & \multicolumn{1}{|l|}{Substraktion} \\ \hline
  \end{tabular}
  \captionof{table}{Table Example 3}
  \label{tab:ex3}
  \endgroup

  \vfill\null\columnbreak

  \begin{minted}{latex}
\begingroup
  \begin{tabular}{ | c | c | }
    \hline
    \multicolumn{1}{|c|}{\emph{ Operator }} & \multicolumn{1}{|c|}{\emph{ Beschreibung}} \\ \hline
    +  & \multicolumn{1}{|l|}{Addition} \\ \hline
    -  & \multicolumn{1}{|l|}{Substraktion} \\ \hline
  \end{tabular}
  \captionof{table}{Table Example 3}
  \label{tab:ex3}
\endgroup
  \end{minted}
\end{multicols}

\begin{multicols}{2}
  \begingroup
  \begin{tabular}{|l||*{2}{c|}}\hline
    \backslashbox{Room}{Date}&\makebox[3em]{Col 1}&\makebox[3em]{Col 2} \\ \hline\hline
    Row 1 &&\\\hline
    Row 2 &&\\\hline
    Row 3 &&\\\hline
  \end{tabular}
  \captionof{table}{Table Example 4}
  \label{tab:ex4}
  \endgroup

  \vfill\null\columnbreak

  \begin{minted}{latex}
\begingroup
\begin{tabular}{|l||*{2}{c|}}\hline
  \backslashbox{Room}{Date}&\makebox[3em]{Col 1}&\makebox[3em]{Col 2} \\ \hline\hline
  Row 1 &&\\\hline
  Row 2 &&\\\hline
  Row 3 &&\\\hline
\end{tabular}
\captionof{table}{Table Example 4}
\label{tab:ex4}
\endgroup
\end{minted}
\end{multicols}

% -----------------------------------------------------------------------------
\chapter{Colors}
\minitoc\newpage

\begin{center}
    \begin{tabular}{l | l}
        Latex Symbol                             & Latex Code                  \\ \hline \hline
        \textcolor{black}{black}                 & \mintinline{latex}{\textcolor{black}{black}} \\
        \textcolor{white}{white}                 & \mintinline{latex}{\textcolor{white}{white}} \\
        \textcolor{red}{red}                     & \mintinline{latex}{\textcolor{red}{red}} \\
        \textcolor{yellow}{yellow}               & \mintinline{latex}{\textcolor{yellow}{yellow}} \\
        \textcolor{lime}{lime}                   & \mintinline{latex}{\textcolor{lime}{lime}} \\
        \textcolor{olive}{olive}                 & \mintinline{latex}{\textcolor{olive}{olive}} \\
        \textcolor{green}{green}                 & \mintinline{latex}{\textcolor{green}{green}} \\
        \textcolor{teal}{teal}                   & \mintinline{latex}{\textcolor{teal}{teal}} \\
        \textcolor{blue}{blue}                   & \mintinline{latex}{\textcolor{blue}{blue}} \\
        \textcolor{HEICyan}{HEICyan}             & \mintinline{latex}{\textcolor{HEICyan}{HEICyan}} \\
        \textcolor{HEIMagenta}{HEIMagenta}       & \mintinline{latex}{\textcolor{HEIMagenta}{HEIMagenta}} \\
        \textcolor{HEIYellow}{HEIYellow}         & \mintinline{latex}{\textcolor{HEIYellow}{HEIYellow}} \\
        \textcolor{HEIGreen}{HEIGreen}           & \mintinline{latex}{\textcolor{HEIGreen}{HEIGreen}} \\
        \textcolor{SPLGreen}{SPLGreen}           & \mintinline{latex}{\textcolor{SPLGreen}{SPLGreen}} \\
        \textcolor{SPLBlue}{SPLBlue}             & \mintinline{latex}{\textcolor{SPLBlue}{SPLBlue}} \\
        \textcolor{SPLPurple}{SPLPurple}         & \mintinline{latex}{\textcolor{SPLPurple}{SPLPurple}} \\
        \textcolor{mGray20}{mGray20}             & \mintinline{latex}{\textcolor{mGray20}{mGray20}} \\
        \textcolor{mGray40}{mGray40}             & \mintinline{latex}{\textcolor{mGray40}{mGray40}} \\
        \textcolor{mGray60}{mGray60}             & \mintinline{latex}{\textcolor{mGray60}{mGray60}} \\
        \textcolor{mGray80}{mGray80}             & \mintinline{latex}{\textcolor{mGray80}{mGray80}} \\
        \textcolor{mWhite}{mWhite}               & \mintinline{latex}{\textcolor{mWhite}{mWhite}} \\
        \textcolor{mBlack}{mBlack}               & \mintinline{latex}{\textcolor{mBlack}{mBlack}} \\
        \textcolor{mPink}{mPink}                 & \mintinline{latex}{\textcolor{mPink}{mPink}} \\
        \textcolor{coloraccent1}{Accent Color 1} & \mintinline{latex}{\textcolor{coloraccent1}{Accent Color 1}} \\
        \textcolor{coloraccent2}{Accent Color 2} & \mintinline{latex}{\textcolor{coloraccent2}{Accent Color 2}} \\
        \textcolor{coloraccent3}{Accent Color 3} & \mintinline{latex}{\textcolor{coloraccent3}{Accent Color 3}} \\
    \end{tabular}
\end{center}

\begin{multicols}{2}
    %\pagecolor{black}  % would make this and all comming pages black
    {
        \color{mGray80}
        Test text \\
    }
    {
        \color{coloraccent1}
        \rule{\linewidth}{1mm}
    }
    \colorbox{coloraccent2}{colorbox}

    \vfill\null\columnbreak

    \begin{minted}{latex}
%\pagecolor{black}  % would make this and all comming pages black
{
\color{mGray80}
Test text \\
}
{
\color{coloraccent1}
\rule{\linewidth}{1mm}
}
\colorbox{coloraccent2}{colorbox}
    \end{minted}
\end{multicols}

% -----------------------------------------------------------------------------
\chapter{Code}
\minitoc\newpage

\section{Inline Code}
\begin{multicols}{2}
  Inline minted \mintinline{cpp}{int x = 0} \\
  Inline code \lstinline{int x = 0}

  \vfill\null\columnbreak

  \begin{minted}{latex}
Inline minted \mintinline{cpp}{int x = 0}
Inline code \lstinline{int x = 0}
  \end{minted}
\end{multicols}

\section{Bloc Code}

\begin{multicols}{2}
  \begin{minted}
[
fontsize=\scriptsize,
linenos
]{python}
import numpy as np
class PiClass:
""" Pi class for getting pi value """
  def __init__():
  """ Returns value of Pi """
    return np.pi
\end{minted}

    \vfill\null\columnbreak

    \begin{minted}{latex}
\begin{minted}
[
fontsize=\scriptsize,
linenos
]{python}
import numpy as np
  class PiClass:
  """ Pi class for getting pi value """
    def __init__():
    """ Returns value of Pi """"
      return np.pi
\end{minte d}
    \end{minted}
\end{multicols}

\begin{multicols}{2}
  \usemintedstyle{monokai}
  \begin{minted}[bgcolor=black!80,]{python}
import numpy as np
class PiClass:
""" Pi class for getting pi value """
  def __init__():
  """ Returns value of Pi """
    return np.pi
  \end{minted}
  \usemintedstyle{\codecolorscheme}
  \vfill\null\columnbreak
  \begin{minted}{latex}
\usemintedstyle{monokai}
\begin{minted}[bgcolor=black!80]{python}
import numpy as np
class PiClass:
""" Pi class for getting pi value """
  def __init__():
  """ Returns value of Pi """
    return np.pi
\end{minte d}
    \end{minted}
\end{multicols}

\begin{multicols}{2}
  \usemintedstyle{bw}
  \begin{minted}
  [  frame=lines,
  framesep=2mm,
  baselinestretch=1.2,
  bgcolor=gray!20,]{python}
import numpy as np
class PiClass:
""" Pi class for getting pi value """
  def __init__():
  """ Returns value of Pi """
    return np.pi
  \end{minted}
  \usemintedstyle{\codecolorscheme}
  \vfill\null\columnbreak

  \begin{minted}{latex}
\usemintedstyle{bw}
\begin{minted}
[frame=lines,
framesep=2mm,
baselinestretch=1.2,
bgcolor=gray!20,]{python}
import numpy as np
class PiClass:
""" Pi class for getting pi value """
  def __init__():
  """ Returns value of Pi """
    return np.pi
\end{minte d}
  \end{minted}
\end{multicols}

\section{Listings}

\begin{multicols}{2}
  \begin{center}
    \inputminted{java}{02-main/listings/HelloWorld.java}
    \captionof{listing}{listing example}
    \label{lst:listing_example}
  \end{center}
  \vfill\null\columnbreak
  \begin{minted}{latex}
\begin{center}
  \inputminted{java}{02-main/listings/HelloWorld.java}
  \captionof{listing}{listing example}
  \label{lst:listing_example}
\end{center}
  \end{minted}

\end{multicols}

% -----------------------------------------------------------------------------
\chapter{MathJax}
\minitoc\newpage

\begin{itemize}
  \item \href{docs.mathjax.org/en/latest/tex.html}{Mathjax Latest Documentation}
  \item \href{https://math.meta.stackexchange.com/questions/5020/mathjax-basic-tutorial-and-quick-reference}{Stack Exchange Basic Tutorial and Quick Reference}
  \item \href{https://oeis.org/wiki/List_of_LaTeX_mathematical_symbols}{List of LaTeX Mathematical Symbols}
  \item \href{https://latex.wikia.org/wiki/List_of_LaTeX_symbols}{List of LaTeX Symbols}
\end{itemize}

\section{Equations}

\subsection{Environments}
The package \mintinline{latex}{amsmath} features 2 different environments \mintinline{latex}{align}, \mintinline{latex}{split}, \mintinline{latex}{multiline}

\subsubsection{align} \newline
Numbered equations aligned at points marked with \mintinline{latex}{\&} usually just before a relation.
\begin{multicols}{2}
  \begin{align}
    a_1& =b_1+c_1\\
    a_2& =b_2+c_2-d_2+e_2
  \end{align}
  \vfill\null\columnbreak
  \begin{minted}{latex}
\begin{align}
  a_1& =b_1+c_1\\
  a_2& =b_2+c_2-d_2+e_2
\end{align}
  \end{minted}
\end{multicols}

\subsubsection{split} \newline
Similar alignment to \mintinline{latex}{align}, but the whole construct fits within equation and is numbered as a unit.
\begin{multicols}{2}
  \begin{equation}
    \begin{split}
      a& =b+c-d\\
       & \quad +e-f\\
       & =g+h\\
       & =i
    \end{split}
  \end{equation}
  \vfill\null\columnbreak
  \begin{minted}{latex}
\begin{equation}
  \begin{split}
    a& =b+c-d\\
     & \quad +e-f\\
     & =g+h\\
     & =i
  \end{split}
\end{equation}
  \end{minted}
\end{multicols}

\subsubsection{multline} \newline
For long expressions taking more than one line, with no specified alignment points.
\begin{multicols}{2}
  \begin{multline}
    a+b+c+d+e+f+g\\
    l+m+n+o+p+q+r
  \end{multline}
  \vfill\null\columnbreak
  \begin{minted}{latex}
\begin{multline}
  a+b+c+d+e+f+g\\
  l+m+n+o+p+q+r
\end{multline}
  \end{minted}
\end{multicols}

\section{Symbols}
Mathjax code need to be places between \$ Symbols
\begin{itemize}
  \item \mintinline{latex}{$...$}: for inline Mathjax
  \item \mintinline{latex}{$$...$$}: for bloc Mathjax
\end{itemize}

\section{Text \& Additions}
\begin{center}
  \begin{tabular}{l | l V{2.7} l | l}
    Symbol                     & Code                                         & Symbol                       & Code                                           \\ \hline \hline
    $\text{normal text}$       & \mintinline{latex}{\text{normal text}}       & $\textbf{bold text}$         & \mintinline{latex}{\textbf{bold text}}         \\
    $\textit{italic text}$     & \mintinline{latex}{\textit{italic text}}     & $\texttt{fixspace text}$     & \mintinline{latex}{\texttt{fixspace text}}     \\
    $Q_1$ $I_{min}$            & \mintinline{latex}{Q_1 I_{min}}              & $x^2$ $x^{20}$               & \mintinline{latex}{x^2 x^{20}}                 \\
    $\acute{x}$                & \mintinline{latex}{\acute{x}}                & $\grave{x}$                  & \mintinline{latex}{\grave{x}}                  \\
    $\dot{x}$                  & \mintinline{latex}{\dot{x}}                  & $\ddot{x}$                   & \mintinline{latex}{\ddot{x}}                   \\
    $\bar{x}$                  & \mintinline{latex}{\bar{x}}                  & $\tilde{x}$                  & \mintinline{latex}{\tilde{x}}                  \\
    $\hat{x}$                  & \mintinline{latex}{\hat{x}}                  & $\check{x}$                  & \mintinline{latex}{\check{x}}                  \\
    $\vec{x}$                  & \mintinline{latex}{\vec{x}}                  & $\breve{x}$                  & \mintinline{latex}{\breve{x}}                  \\
    $\overset{over}{X}$        & \mintinline{latex}{\overset{over}{X}}        & $\underset{under}{X}$        & \mintinline{latex}{\underset{under}{X}}        \\
    $\overline{xxx}$           & \mintinline{latex}{\overline{xxx}}           & $\underline{xxx}$            & \mintinline{latex}{\underline{xxx}}            \\
    $\overbrace{xxx}$          & \mintinline{latex}{\overbrace{xxx}}          & $\underbrace{xxx}$           & \mintinline{latex}{\underbrace{xxx}}           \\
    $\overleftarrow{xxx}$      & \mintinline{latex}{\overleftarrow{xxx}}      & $\overrightarrow{xxx}$       & \mintinline{latex}{\overrightarrow{xxx}}       \\
    $\xleftarrow[under]{over}$ & \mintinline{latex}{\xleftarrow[under]{over}} & $\xrightarrow[under]{over}$  & \mintinline{latex}{\xrightarrow[under]{over}}  \\
  \end{tabular}
\end{center}

\section{Spaces}
\begin{center}
  \begin{tabular}{l | l V{2.7} l | l V{2.7} l | l}
    Symbol              & Code                                  & Symbol              & Code                                   & Symbol               & Code                 \\ \hline \hline
    $a \mspace{3mu} b$  & \mintinline{latex}{a \mspace{3mu} b}  & $a \, b$            & \mintinline{latex}{a \, b}             & $a \thinspace b$     & \mintinline{latex}{a \thinspace b}     \\
    $a \mspace{4mu} b$  & \mintinline{latex}{a \mspace{4mu} b}  & $a \:  b$           & \mintinline{latex}{a \: b}             & $a \medspace b$      & \mintinline{latex}{a \medspace b}      \\
    $a \mspace{5mu} b$  & \mintinline{latex}{a \mspace{5mu} b}  & $a \; b$            & \mintinline{latex}{a \; b}             & $a \thickspace b$    & \mintinline{latex}{a \thickspace b}    \\
    $a \mspace{6mu} b$  & \mintinline{latex}{a \mspace{6mu} b}  & $a \ b$             & \mintinline{latex}{a \ b}              & -                    & -                    \\
    $a \mspace{18mu} b$ & \mintinline{latex}{a \mspace{18mu} b} & $a \quad b$         & \mintinline{latex}{a \quad b}          & -                    & -                    \\
    $a \mspace{36mu} b$ & \mintinline{latex}{a \mspace{36mu} b} & $a \qquad b$        & \mintinline{latex}{a \qquad b}         & -                    & -                    \\
    $a \mspace{-3mu} b$ & \mintinline{latex}{a \mspace{-3mu} b} & $a \! b$            & \mintinline{latex}{a \, b}             & $a \negthinspace b$  & \mintinline{latex}{a \negthinspace b}  \\
    $a \mspace{-4mu} b$ & \mintinline{latex}{a \mspace{-4mu} b} & -                   & -                    & $a \negmedspace b$   & \mintinline{latex}{a \negmedspace b}   \\
    $a \mspace{-5mu} b$ & \mintinline{latex}{a \mspace{-5mu} b} & -                   & -                    & $a \negthickspace b$ & \mintinline{latex}{a \negthickspace b} \\
  \end{tabular}
\end{center}

\section{Quotes}
\begin{center}
  \begin{tabular}{l | l | l}
    Symbol       & Code                     & Alt Code \\ \hline \hline
    $\q$         & \mintinline{latex}{\q}   & \mintinline{latex}{\textnormal{\textquotesingle}} \\
    $\ql$        & \mintinline{latex}{\ql}  & \mintinline{latex}{\textnormal{\textquoteleft}} \\
    $\qr$        & \mintinline{latex}{\qr}  & \mintinline{latex}{\textnormal{\textquoteright}} \\
    $\qq$        & \mintinline{latex}{\qq}  & \mintinline{latex}{\textnormal{\textquotedblright}} \\
    $\qql$       & \mintinline{latex}{\qql} & \mintinline{latex}{\textnormal{\textquotedblleft}} \\
    $\qqr$       & \mintinline{latex}{\qqr} & \mintinline{latex}{\textnormal{\textquotedblright}} \\
  \end{tabular}
\end{center}

\section{Greek Letters}
\begin{center}
  \begin{tabular}{l | l}
    Symbol                                & Code \\ \hline \hline
    $A \alpha$                            & \mintinline{latex}{A \alpha} \\
    $B \beta$                             & \mintinline{latex}{B \beta} \\
    $\Gamma \gamma$                       & \mintinline{latex}{\Gamma \gamma} \\
    $\Delta \delta$                       & \mintinline{latex}{\Delta \delta} \\
    $E \epsilon \varepsilon $             & \mintinline{latex}{E \epsilon \varepsilon} \\
    $Z \zeta$                             & \mintinline{latex}{Z \zeta} \\
    $H \eta$                              & \mintinline{latex}{H \eta} \\
    $\Theta \theta \vartheta$             & \mintinline{latex}{\Theta \theta \vartheta} \\
    $I \iota$                             & \mintinline{latex}{I\iota} \\
    $K \kappa$                            & \mintinline{latex}{K \kappa} \\
    $\Lambda \lambda$                     & \mintinline{latex}{\Lambda \lambda} \\
    $M \mu$                               & \mintinline{latex}{M \mu} \\
    $N \nu$                               & \mintinline{latex}{N \nu} \\
    $\Xi \xi$                             & \mintinline{latex}{\Xi \xi} \\
    $O \ominus$                           & \mintinline{latex}{O \ominus} \\
    $\Pi \pi \varpi$                      & \mintinline{latex}{\Pi \pi \varpi} \\
    $P \rho$                              & \mintinline{latex}{P \rho} \\
    $\Sigma \sigma \varsigma$             & \mintinline{latex}{\Sigma \sigma \varsigma} \\
    $T \tau$                              & \mintinline{latex}{T \tau} \\
    $\Upsilon \upsilon$                   & \mintinline{latex}{\Upsilon \upsilon} \\
    $\Phi \phi \varphi$                   & \mintinline{latex}{\Phi \phi \varphi} \\
    $X \chi$                              & \mintinline{latex}{X \chi} \\
    $\Psi \psi$                           & \mintinline{latex}{\Psi \psi} \\
    $\Omega \omega$                       & \mintinline{latex}{\Omega \omega} \\
  \end{tabular}
\end{center}

\section{Symbols}
\subsection{Relational Operators}

\begin{center}
  \begin{tabular}{l | l V{2.7} l | l}
    Symbol                                & Code                            & Symbol                                & Code \\ \hline \hline
    $<$                                   & \mintinline{latex}{<}           & $>$                                   & \mintinline{latex}{>} \\
    $\nless$                              & \mintinline{latex}{\nless}      & $\ngtr$                               & \mintinline{latex}{\ngtr} \\
    $\leq$                                & \mintinline{latex}{\leq}        & $\geq$                                & \mintinline{latex}{\geq} \\
    $\leqslant$                           & \mintinline{latex}{\leqslant}   & $\geqslant$                           & \mintinline{latex}{\geqslant} \\
    $\nleq$                               & \mintinline{latex}{\nleq}       & $\ngeq$                               & \mintinline{latex}{\ngeq} \\
    $\nleqslant$                          & \mintinline{latex}{\nleqslant}  & $\ngeqslant$                          & \mintinline{latex}{\ngeqslant} \\
    $\ll$                                 & \mintinline{latex}{\ll}         & $\gg$                                 & \mintinline{latex}{\gg} \\
    $\lll$                                & \mintinline{latex}{\lll}        & $\ggg$                                & \mintinline{latex}{\ggg} \\
    $\subset$                             & \mintinline{latex}{\subset}     & $\supset$                             & \mintinline{latex}{\supset} \\
    $\not\subset$                         & \mintinline{latex}{\not\subset} & $\not\supset$                         & \mintinline{latex}{\not\supset} \\
    $\subseteq$                           & \mintinline{latex}{\subseteq}   & $\supseteq$                           & \mintinline{latex}{\supseteq} \\
    $\nsubseteq$                          & \mintinline{latex}{\nsubseteq}  & $\nsupseteq$                          & \mintinline{latex}{\nsupseteq} \\
  \end{tabular}

  \begin{tabular}{l | l}
    Symbol                                & Code \\ \hline \hline
    $=$                                   & \mintinline{latex}{=} \\
    $\equiv$                              & \mintinline{latex}{\equiv} \\
    $\approx$                             & \mintinline{latex}{\approx} \\
    $\cong$                               & \mintinline{latex}{\cong} \\
    $\simeq$                              & \mintinline{latex}{\simeq} \\
    $\sim$                                & \mintinline{latex}{\sim} \\
    $\propto$                             & \mintinline{latex}{\propto} \\
    $\neq \ne$                            & \mintinline{latex}{\neq \ne} \\
  \end{tabular}
\end{center}

\subsection{Arrows}
\begin{center}
  \begin{tabular}{l | l V{2.7} l | l}
     Symbol                               &  Code                               &  Symbol             &  Code \\ \hline \hline
    $\rightarrow \to$                     & \mintinline{latex}{\rightarrow \to} & $\leftarrow \gets$  & \mintinline{latex}{\leftarrow \gets} \\
    $\Rightarrow$                         & \mintinline{latex}{\Rightarrow}     & $\Leftarrow$        & \mintinline{latex}{\Leftarrow} \\
    $\longrightarrow$                     & \mintinline{latex}{\longrightarrow} & $\longrightarrow$   & \mintinline{latex}{\longrightarrow} \\
    $\Longrightarrow$                     & \mintinline{latex}{\Longrightarrow} & $\Longrightarrow$   & \mintinline{latex}{\Longrightarrow} \\
    $\mapsto$                             & \mintinline{latex}{\mapsto}         & -                   & - \\
    $\longmapsto$                         & \mintinline{latex}{\longmapsto}     & -                   & - \\
    $\uparrow$                            & \mintinline{latex}{\uparrow}        & $\Uparrow$          & \mintinline{latex}{\Uparrow} \\
    $\downarrow$                          & \mintinline{latex}{\downarrow}      & $\Downarrow$        & \mintinline{latex}{\Downarrow} \\
    $\updownarrow$                        & \mintinline{latex}{\updownarrow}    & $\Updownarrow$      & \mintinline{latex}{\Updownarrow} \\
  \end{tabular}
\end{center}

\subsection{Boolsche Algebra Symbols}
\begin{center}
  \begin{tabular}{l V{2.7} l | l}
    Operator        & Symbol                                 & Code \\ \hline \hline
    NEGATE          & $\neg \; \overline{x} \; \bar{x}$      & \mintinline{latex}{\neg \overline{x} \bar{x}} \\
    AND             & $\bigwedge \; \wedge \; * \; \And$     & \mintinline{latex}{\bigwedge \wedge * \And} \\
    OR              & $\bigvee \; \vee \; + \; \mid$         & \mintinline{latex}{\bigvee \vee + \mid} \\
    XOR             & $\oplus$                               & \mintinline{latex}{\oplus} \\
  \end{tabular}
\end{center}

\subsection{Other Symbols}
\begin{center}
  \begin{tabular}{l | l V{2.7} l | l V{2.7} l | l}
    Symbol          & Code                              & Symbol               & Code                                   & Symbol            & Code \\ \hline \hline
    $\#$            & \mintinline{latex}{\#}            & $\sharp$             & \mintinline{latex}{\sharp}             & $\S$              & \mintinline{latex}{\S} \\
    $\lozenge$      & \mintinline{latex}{\lozenge}      & $\blacklozenge$      & \mintinline{latex}{\blacklozenge}      & $\infty$          & \mintinline{latex}{\infty} \\
    $\square$       & \mintinline{latex}{\square}       & $\blacksquare$       & \mintinline{latex}{\blacksquare}       & $\spadesuit$      & \mintinline{latex}{\spadesuit} \\
    $\triangle$     & \mintinline{latex}{\triangle}     & $\blacktriangle$     & \mintinline{latex}{\blacktriangle}     & $\clubsuit$       & \mintinline{latex}{\clubsuit} \\
    $\triangledown$ & \mintinline{latex}{\triangledown} & $\blacktriangledown$ & \mintinline{latex}{\blacktriangledown} & $\heartsuit$      & \mintinline{latex}{\heartsuit} \\
    $\diagup$       & \mintinline{latex}{\diagup}       & $\diagdown$          & \mintinline{latex}{\diagdown}          & $\diamondsuit$    & \mintinline{latex}{\diamondsuit} \\
    $\varnothing$   & \mintinline{latex}{\varnothing}   & $\emptyset$          & \mintinline{latex}{\emptyset}          & $\angle$          & \mintinline{latex}{\angle} \\
    $\square$       & \mintinline{latex}{\square}       & $\surd$              & \mintinline{latex}{\surd}              & $\measuredangle$  & \mintinline{latex}{\measuredangle} \\
  \end{tabular}
\end{center}

\section{Math Symbols}
\subsection{Trigonometry}
\begin{center}
  \begin{tabular}{l | l V{2.7} l | l V{2.7} l | l}
    Symbol       & Code                         & Symbol       & Code                         & Symbol       & Code \\ \hline \hline
    $\sin$       & \mintinline{latex}{\sin}     & $\arcsin$    & \mintinline{latex}{\arcsin}  & $\sinh$      & \mintinline{latex}{\sinh} \\
    $\cos$       & \mintinline{latex}{\cos}     & $\arccos$    & \mintinline{latex}{\arccos}  & $\cosh$      & \mintinline{latex}{\cosh} \\
    $\tan$       & \mintinline{latex}{\tan}     & $\arctan$    & \mintinline{latex}{\arctan}  & $\tanh$      & \mintinline{latex}{\tanh} \\
  \end{tabular}
\end{center}

\subsection{Prefix Operators}
\begin{center}
  \begin{tabular}{l | l V{2.7} l | l V{2.7} l | l}
    Symbol       & Code                         & Symbol       & Code                          & Symbol       & Code   \\ \hline \hline
    $\int$       & \mintinline{latex}{\int}     & $\oint$      & \mintinline{latex}{\oint}     & $\sum$       & \mintinline{latex}{\sum}       \\
    $\prod$      & \mintinline{latex}{\prod}    & $\coprod$    & \mintinline{latex}{\coprod}   & -            & -                              \\
    $\bigodot$   & \mintinline{latex}{\bigodot} & $\bigoplus$  & \mintinline{latex}{\bigoplus} & $\bigotimes$ & \mintinline{latex}{\bigotimes} \\
    $\bigcap$    & \mintinline{latex}{\bigcap}  & $\bigcup$    & \mintinline{latex}{\bigcup}   & $\bigsqcup$  & \mintinline{latex}{\bigsqcup}  \\
    $\bigvee$    & \mintinline{latex}{\bigvee}  & $\bigwedge$  & \mintinline{latex}{\bigwedge} & -            & -                              \\
  \end{tabular}
\end{center}

% -----------------------------------------------------------------------------
\chapter{Fancy Additions}
\minitoc\newpage

\section{Comments}
Comments can be activated and deactivated with the options packet in the file \lstinline{01-settings/metadata.tex}

\begin{itemize}
  \item \mintinline{latex}{\todo[author=tschinz]{Todo Comment with Authorname}} \todo[author=tschinz]{Todo Comment with Authorname}
  \item \mintinline{latex}{\todo{Todo Comment}} \todo{Todo Comment}
  \item \mintinline{latex}{\unsure{Unsure Comment}} \unsure{Unsure Comment}
  \item \mintinline{latex}{\info{Info Comment}} \info{Info Comment}
  \item \mintinline{latex}{\critical{Critical Comment}} \critical{Critical Comment}
  \item \mintinline{latex}{\change{Change Comment}} \change{Change Comment}
  \item \mintinline{latex}{\question{Question Comment}} \question{Question Comment}
  \item \mintinline{latex}{\improvement{Improvement Comment}} \improvement{Improvement Comment}
\end{itemize}

\section{Boxes}
The package awesomebox allows for some nice looking boxes.

\begin{multicols}{2}
  \notebox{Lorem Ipsum ...}
  \vfill\null\columnbreak
  \begin{minted}{latex}
\notebox{Lorem Ipsum ...}
  \end{minted}
\end{multicols}

\begin{multicols}{2}
  \tipbox{Lorem Ipsum ...}
  \vfill\null\columnbreak
  \begin{minted}{latex}
\tipbox{Lorem Ipsum ...}
  \end{minted}
\end{multicols}

\begin{multicols}{2}
  \warningbox{Lorem Ipsum ...}
  \vfill\null\columnbreak
  \begin{minted}{latex}
\warningbox{Lorem Ipsum ...}
  \end{minted}
\end{multicols}

\begin{multicols}{2}
  \cautionbox{Lorem Ipsum ...}
  \vfill\null\columnbreak
  \begin{minted}{latex}
\cautionbox{Lorem Ipsum ...}
  \end{minted}
\end{multicols}

\begin{multicols}{2}
  \importantbox{Lorem Ipsum ...}
  \vfill\null\columnbreak
  \begin{minted}{latex}
\importantbox{Lorem Ipsum ...}
  \end{minted}
\end{multicols}

\begin{multicols}{2}
  \awesomebox[violet]{2pt}{\faRocket}{violet}{Lorem Ipsum ...}
  \vfill\null\columnbreak
  \begin{minted}{latex}
\awesomebox[violet]{2pt}{\faRocket}{violet}{Lorem Ipsum ...}
  \end{minted}
\end{multicols}

\section{Package Pifont Special Characters}
Insert with \mintinline{latex}{\ding{0}}
\begin{center}
  \begin{tabular}{l  c  c  c  c  c  c  c  c}
    Index & 0        & 1          & 2          & 3          & 4          & 5          & 6          & 7 \\ \hline\hline
    0   & \ding{0}   & \ding{1}   & \ding{2}   & \ding{3}   & \ding{4}   & \ding{5}   & \ding{6}   & \ding{7} \\
    8   & \ding{8}   & \ding{9}   & \ding{10}  & \ding{11}  & \ding{12}  & \ding{13}  & \ding{14}  & \ding{15} \\
    16  & \ding{16}  & \ding{17}  & \ding{18}  & \ding{19}  & \ding{20}  & \ding{21}  & \ding{22}  & \ding{23} \\
    24  & \ding{24}  & \ding{25}  & \ding{26}  & \ding{27}  & \ding{28}  & \ding{29}  & \ding{30}  & \ding{31} \\
    32  & \ding{32}  & \ding{33}  & \ding{34}  & \ding{35}  & \ding{36}  & \ding{37}  & \ding{38}  & \ding{39} \\
    40  & \ding{40}  & \ding{41}  & \ding{42}  & \ding{43}  & \ding{44}  & \ding{45}  & \ding{46}  & \ding{47} \\
    48  & \ding{48}  & \ding{49}  & \ding{50}  & \ding{51}  & \ding{52}  & \ding{53}  & \ding{54}  & \ding{55} \\
    56  & \ding{56}  & \ding{57}  & \ding{58}  & \ding{59}  & \ding{60}  & \ding{61}  & \ding{62}  & \ding{63} \\
    64  & \ding{64}  & \ding{65}  & \ding{66}  & \ding{67}  & \ding{68}  & \ding{69}  & \ding{70}  & \ding{71} \\
    72  & \ding{72}  & \ding{73}  & \ding{74}  & \ding{75}  & \ding{76}  & \ding{77}  & \ding{78}  & \ding{79} \\
    80  & \ding{80}  & \ding{81}  & \ding{82}  & \ding{83}  & \ding{84}  & \ding{85}  & \ding{86}  & \ding{87} \\
    88  & \ding{88}  & \ding{89}  & \ding{90}  & \ding{91}  & \ding{92}  & \ding{93}  & \ding{94}  & \ding{95} \\
    96  & \ding{96}  & \ding{97}  & \ding{98}  & \ding{99}  & \ding{100} & \ding{101} & \ding{102} & \ding{103} \\
    104 & \ding{104} & \ding{105} & \ding{106} & \ding{107} & \ding{108} & \ding{109} & \ding{110} & \ding{111} \\
    112 & \ding{112} & \ding{113} & \ding{114} & \ding{115} & \ding{116} & \ding{117} & \ding{118} & \ding{119} \\
    120 & \ding{120} & \ding{121} & \ding{122} & \ding{123} & \ding{124} & \ding{125} & \ding{126} & \ding{127} \\
    128 & \ding{128} & \ding{129} & \ding{130} & \ding{131} & \ding{132} & \ding{133} & \ding{134} & \ding{135} \\
    136 & \ding{136} & \ding{137} & \ding{138} & \ding{139} & \ding{140} & \ding{141} & \ding{142} & \ding{143} \\
    144 & \ding{144} & \ding{145} & \ding{146} & \ding{147} & \ding{148} & \ding{149} & \ding{150} & \ding{151} \\
    152 & \ding{152} & \ding{153} & \ding{154} & \ding{155} & \ding{156} & \ding{157} & \ding{158} & \ding{159} \\
    160 & \ding{160} & \ding{161} & \ding{162} & \ding{163} & \ding{164} & \ding{165} & \ding{166} & \ding{167} \\
    168 & \ding{168} & \ding{169} & \ding{170} & \ding{171} & \ding{172} & \ding{173} & \ding{174} & \ding{175} \\
    176 & \ding{176} & \ding{177} & \ding{178} & \ding{179} & \ding{180} & \ding{181} & \ding{182} & \ding{183} \\
    184 & \ding{184} & \ding{185} & \ding{186} & \ding{187} & \ding{188} & \ding{189} & \ding{190} & \ding{191} \\
    192 & \ding{192} & \ding{193} & \ding{194} & \ding{195} & \ding{196} & \ding{197} & \ding{198} & \ding{199} \\
    200 & \ding{200} & \ding{201} & \ding{202} & \ding{203} & \ding{204} & \ding{205} & \ding{206} & \ding{207} \\
    208 & \ding{208} & \ding{209} & \ding{210} & \ding{211} & \ding{212} & \ding{213} & \ding{214} & \ding{215} \\
    216 & \ding{216} & \ding{217} & \ding{218} & \ding{219} & \ding{220} & \ding{221} & \ding{222} & \ding{223} \\
    224 & \ding{224} & \ding{225} & \ding{226} & \ding{227} & \ding{228} & \ding{229} & \ding{230} & \ding{231} \\
    232 & \ding{232} & \ding{233} & \ding{234} & \ding{235} & \ding{236} & \ding{237} & \ding{238} & \ding{239} \\
    230 & \ding{240} & \ding{241} & \ding{242} & \ding{243} & \ding{244} & \ding{245} & \ding{246} & \ding{247} \\
    248 & \ding{248} & \ding{249} & \ding{250} & \ding{251} & \ding{252} & \ding{253} & \ding{254} & \ding{255} \\
  \end{tabular}
\end{center}

\section{Awesome Font}
Symbols from \href{https://fontawesome.com/}{Awesomefont} can be added by \mintinline{latex}{\fa<IconName>}

\begin{multicols}{2}
  \faStar \\
  \faPlusCircle \\
  \faMinusCircle \\
  \faCode \\
  \faRocket \\
  \vfill\null\columnbreak
  \begin{minted}{latex}
\faStar
\faPlusCircle
\faMinusCircle
\faCode
\faRocket
  \end{minted}
\end{multicols}

% -----------------------------------------------------------------------------
\chapter{References}
\minitoc\newpage

\section{Acronyms}
\begin{center}
  \begin{tabular}{l | l V{2.7} p{6cm} }
    Latex Output  & Latex Code                         & Description \\ \hline \hline
    \acrlong{ar}  & \mintinline{latex}{\acrlong{ar}}   & Displays the phrase which the acronyms stands for. Put the label of the acronym inside the braces. \\ \hline
    \acrshort{ar} & \mintinline{latex}{\acrshort{ar}}  & Prints the acronym whose label is passed as parameter. \\ \hline
    \acrfull{ar}  & \mintinline{latex}{\acrfull{ar}}   & Prints both, the acronym and its definition. \\
  \end{tabular}
\end{center}

\section{Glossary}
\begin{center}
  \begin{tabular}{l | l V{2.7} p{8cm} }
    Latex Output   & Latex Code & Description \\ \hline \hline
    \gls{scrum}   & \mintinline{latex}{\gls{scrum}}    & To print the term, lowercase. For example, \gls{scrum} prints mathematics when used. \\ \hline
    \Gls{scrum}   & \mintinline{latex}{\Gls{scrum}}    & The same as \gls but the first letter will be printed in uppercase. Example: \Gls{scrum} prints Mathematics \\ \hline
    \glspl{scrum} & \mintinline{latex}{\glspl{scrum}}  & The same as \gls but the term is put in its plural form. For instance, \glspl{scrum} will write formulas in your final document. \\ \hline
    \Glspl{scrum} & \mintinline{latex}{\Glspl{scrum}}  & The same as \Gls but the term is put in its plural form. For example, \Glspl{scrum} renders as Formulas. \\ \hline
  \end{tabular}
\end{center}

\section{Bibliography}
\begin{center}
  \begin{tabular}{l | l V{2.7} p{4cm} }
    Output                              & Code                                                    & Description \\ \hline \hline
    \cite{stateoftheArt}                & \mintinline{latex}{\cite{stateoftheArt}}                & Default Citations \\
    \cite[before][after]{stateoftheArt} & \mintinline{latex}{\cite[before][after]{stateoftheArt}} & Citation with additions before and after \\
  \end{tabular}
\end{center}

% -----------------------------------------------------------------------------
% Back matter
% -----------------------------------------------------------------------------
\backmatter

% Bibliography
\opt{bib}{
  \newpage
  \printbibliography[heading=bibintoc]
}

% Glossary
\opt{acr}{
  \newpage
  \printglossary[type=\acronymtype]
}
\opt{glo}{
  \printglossary
}

\end{document}

